% META: {"id": "TEXP-ARI-CO-011", "chapitre": "TEXP-ARITHMETIQUE", "type_objet": "corrige", "exercice_ref": "TEXP-ARI-EX-011", "capacites_codes": ["C2"], "sources_inspiration": [], "status": "generated"}
% BEGIN-VERIFY
% import math
% assert math.gcd(6, 15) == 3
% assert not any((6 * u) % 15 == 1 for u in range(15))
% for b in range(15):
%     a_des_solutions = any((6 * x) % 15 == b for x in range(15))
%     assert a_des_solutions == (b % 3 == 0)
% assert [x for x in range(15) if (6 * x) % 15 == 9] == [4, 9, 14]
% assert (2 * 3) % 5 == 1
% END-VERIFY
\begin{corrige}{TEXP-ARI-EX-011}
\textbf{1.} $\mathrm{PGCD}(6\,;\,15)=3\neq1$ : $6$ n'est pas inversible
modulo $15$, car un inverse $u$ vérifierait $6u-15v=1$, impossible puisque $3$
divise le membre de gauche.

\textbf{2.} Si $6x\equiv b\ [15]$, alors $b=6x-15k$ pour un entier $k$, et $3$
divise $6x-15k$ : donc $3$ divise $b$.

\textbf{3.} $3$ divise $9$ : il y a des solutions. En divisant tout par $3$,
la congruence équivaut à $2x\equiv3\ [5]$. Comme $2\times3=6\equiv1\ [5]$,
l'inverse de $2$ modulo $5$ est $3$, d'où $x\equiv9\equiv4\ [5]$. Modulo $15$,
cela donne trois classes : $x\in\{4,9,14\}$.
\end{corrige}
